\section{十四、团队介绍与分工}

\subsection{1.团队构成与分工}

本项目由 4 名成员组成，覆盖算法、后端、前端与数据处理等方向，形成分工明确、协同高效的研发团队：由项目负责人统筹全局并主导核心技术攻关，后端成员承担服务端开发与测试，前端与数据成员负责界面实现与数据集构建。

\begin{table}[htbp]
\centering
\caption{团队成员分工与贡献}
\begin{tabular}{|p{1.9cm}|p{2.6cm}|p{6.6cm}|p{3.0cm}|}
\hline
\textbf{姓名} & \textbf{角色} & \textbf{主要工作} & \textbf{贡献说明} \\
\hline
汤福连 & 项目负责人 & 全局统筹；系统架构设计；模型训练（BIT，ResNet18 骨干）；后端 API 核心开发；前端总体设计；AI 与 GIS 集成；测试与安全加固；云服务器部署上线 & 项目核心，独立完成模型、后端、前端主体与部署上线的绝大部分工作 \\
\hline
刘芮仲 & 后端开发 & 后端部分接口开发；测试用例编写；文档与资料整理 & 参与后端开发与测试，辅助主干工作 \\
\hline
张珏枫 & 后端开发 & 后端辅助接口开发；数据库与部署运维协助；问题排查 & 参与后端辅助与部署协作 \\
\hline
于欣欣 & 前端与数据标注 & 前端页面辅助开发；高德地图 API 集成；自建黑土耕地数据集像素级标注 & 参与前端与数据标注工作 \\
\hline
\end{tabular}
\end{table}

从分工结构看，团队以“算法—后端—前端—数据”为主线组织协作：算法与架构由项目负责人统一把关，保证技术路线的一致性；后端成员在统一接口规范下并行开发，避免相互阻塞；前端与数据成员则分别对接系统界面与数据集构建，形成从数据到应用的完整链条。这一结构在项目周期紧、任务跨度大的条件下，有效保证了研发进度与成果质量。

\subsection{2.协作机制}

团队在研发过程中建立了规范、可追溯的协作机制，具体包括：

\begin{itemize}
    \item \textbf{版本管理}：全程使用 Git 进行代码管理，提交采用语义化格式（feat/fix/chore），关键版本打标签，全部变更可回溯、可审计。
    \item \textbf{接口先行}：前后端以 OpenAPI 接口文档为契约并行开发，前端按接口约定独立推进，降低联调耦合与等待成本。
    \item \textbf{数据双人复核}：自建黑土耕地数据集采用“独立标注 + 交叉验证”流程，由两名成员分别标注并交叉核对，降低标注噪声对模型精度的影响。
    \item \textbf{环境统一}：开发、测试与生产环境通过容器化与统一配置模板保持一致，避免“本地可运行、上线即报错”的常见问题。
    \item \textbf{工具提效}：研发过程中引入 AI 编码辅助工具参与需求梳理、代码审查与安全加固，提升迭代效率与代码质量。
\end{itemize}

上述机制使团队在有限人力条件下完成了从模型训练、系统开发到安全加固、部署上线的全流程工作，也为成果的后续维护与推广保留了良好的工程基础。
